\documentclass[11pt]{amsart}
\usepackage[T1]{fontenc}
\usepackage{lmodern}
\usepackage{amsmath,amssymb,amsthm,mathtools}
\usepackage{booktabs,longtable,array,microtype}
\usepackage{xcolor}
\usepackage[colorlinks=true,linkcolor=blue!55!black,citecolor=blue!55!black,urlcolor=blue!55!black]{hyperref}

\newtheorem{theorem}{Theorem}[section]
\newtheorem{lemma}[theorem]{Lemma}
\newtheorem{proposition}[theorem]{Proposition}
\newtheorem{corollary}[theorem]{Corollary}
\theoremstyle{definition}
\newtheorem{definition}[theorem]{Definition}
\newtheorem{remark}[theorem]{Remark}

\DeclareMathOperator{\Var}{Var}
\newcommand{\Cpath}[1]{\mathbf C_{P_{#1}}}
\newcommand{\one}{\mathbf 1}

\title[The sharp variation constant on $P_{11}$]
{The sharp variation constant for the centered Hardy--Littlewood maximal operator on $P_{11}$}
\author{Zijian Zeng}
\address{Institute of Computer Science and Digital Innovation, UCSI University, Kuala Lumpur 56000, Malaysia}
\email{zijianzeng@foxmail.com}
\email{1002266693@ucsiuniversity.edu.my}
\date{August 30, 2026}
\hypersetup{
  pdftitle={The sharp variation constant for the centered Hardy--Littlewood maximal operator on P11},
  pdfauthor={Zijian Zeng}
}

\begin{document}

\begin{abstract}
For the path graph $P_n$, let $\mathbf C_{P_n}$ be the optimal constant in
the $1$-variation inequality for the centered graph Hardy--Littlewood maximal
operator.  It has been conjectured that $\mathbf C_{P_n}=1-1/n$ for every
$n\geq3$.  We prove the next finite case
\[
\mathbf C_{P_{11}}=\frac{10}{11}.
\]
The upper bound is reduced to a finite family of signed edge-flux
certificates.  An exact, independently checked enumeration covers
$5{,}874{,}849$ chambers and uses only rational arithmetic at the
verification stage.  The matching lower bound is attained by a point mass at
an endpoint.  No assertion is made for arbitrary $n$.
\end{abstract}

\maketitle

\section{Introduction}

Let $G=(V,E)$ be a finite connected undirected graph.  The centered graph
Hardy--Littlewood maximal operator averages the absolute value of a function
over graph-distance balls.  Liu and Xue studied sharp variation constants for
this operator on finite graphs~\cite{LiuXue2021}.  González-Riquelme, Kovač,
and Madrid subsequently determined further sharp constants and posed the
path-graph conjecture~\cite[Conjecture~5]{GKM2026}:
\begin{equation}\label{eq:path-conjecture}
  \Cpath n=1-\frac1n\qquad(n\geq3).
\end{equation}
Their paper proves the case $P_4$ by finite linear programming and reports
numerical evidence through $n=10$~\cite[Section~3.2 and Conjecture~5]{GKM2026}.

A directly adjacent infinite-lattice development is Raghava's conditional
Abel-summation edge-flux reduction for centered maximal variation on
$\mathbb Z$~\cite[Abstract and Sections~9--10]{Raghava2026}.  It leaves a
residual packing principle conjectural and supplies no $P_{11}$ certificate.
Its peak-based conditional reduction on the infinite lattice is distinct
from the complete finite rational chamber certificate constructed here.

The purpose of this paper is to determine the next finite parameter.

\begin{theorem}\label{thm:main}
For every real-valued function $f$ on the eleven-vertex path,
\[
  \Var(M_{P_{11}}f)\leq\frac{10}{11}\Var(f).
\]
The constant $10/11$ is optimal and is attained by a point mass at either
endpoint.  Consequently,
\[
  \boxed{\Cpath{11}=\frac{10}{11}}.
\]
\end{theorem}

The proof separates a general, human-checkable chamber lemma from a finite
exact certificate.  After quotienting vertical translations, a function is
represented by its ten edge differences.  The variation of the maximal
function is the maximum of $2^{10}$ signed linear expressions.  Once one
active ball is fixed at each positive coefficient, negative coefficients can
be bounded by convex combinations of centered ball averages.  The resulting
signed measure has total mass zero, and its cumulative masses are precisely
edge-flux coefficients.  Bounding their magnitudes by $10/11$ proves the
desired chamber inequality.

The finite certificate is not a floating-point optimization result.  Its
independent verifier reconstructs rational vertex measures, checks all mass
constraints, and checks every cumulative edge coefficient exactly.  The
complete conjecture~\eqref{eq:path-conjecture} remains open.

\section{Definitions and exact reductions}

Write
\[
  V_n=\{1,\ldots,n\},\qquad
  E_n=\bigl\{\{i,i+1\}:1\leq i<n\bigr\}
\]
for the path $P_n$.  Its graph distance is $d(i,j)=|i-j|$, and
\[
  B_n(i,r)=\{j\in V_n:|i-j|\leq r\}.
\]
Only the integer radii
$0\leq r\leq\max\{i-1,n-i\}$ give distinct balls.  For
$f:V_n\to\mathbb R$, define
\begin{equation}\label{eq:maximal-def}
  (M_{P_n}f)(i)=
  \max_{r\geq0}\frac1{|B_n(i,r)|}
  \sum_{j\in B_n(i,r)}|f(j)|
\end{equation}
and
\begin{equation}\label{eq:variation-def}
  \Var(f)=\sum_{i=1}^{n-1}|f(i+1)-f(i)|.
\end{equation}
The optimal constant is equivalently
\[
  \Cpath n=
  \sup_{\Var(f)>0}\frac{\Var(M_{P_n}f)}{\Var(f)}.
\]
Constant inputs cause no ambiguity: both variations vanish.

For arbitrary real $z$, let $\mathcal A_n z$ denote the same maximum of
centered averages as in~\eqref{eq:maximal-def}, but without putting absolute
values around the entries of $z$.

\begin{lemma}[Translation quotient]\label{lem:quotient}
For $n\geq2$,
\begin{equation}\label{eq:quotient}
  \Cpath n=
  \max_{\substack{z(1)=0\\ \Var(z)=1}}
  \Var(\mathcal A_n z).
\end{equation}
The maximum on the right is attained.
\end{lemma}

\begin{proof}
Since $M_{P_n}f=M_{P_n}|f|$ and
$\Var(|f|)\leq\Var(f)$, it is enough to optimize over nonnegative inputs
$x$.  For such inputs, $M_{P_n}x=\mathcal A_nx$.  Every centered average
commutes with addition of a constant, and therefore
\[
  \mathcal A_n(z+c\one)=\mathcal A_nz+c\one.
\]
Thus a nonnegative $x$ can be replaced by $z=x-x(1)\one$ without changing
either variation.  Conversely, if $z(1)=0$, then choosing
$c\geq-\min_i z(i)$ makes $x=z+c\one$ nonnegative and again preserves both
variations.  Positive homogeneity gives the normalization $\Var(z)=1$.

If $d_k=z(k+1)-z(k)$, then $z(1)=0$ and $\sum_k|d_k|=1$ imply
$|z(i)|\leq1$.  Hence the normalized feasible set is compact, and continuity
of the finite maximum of linear functions proves attainment.
\end{proof}

For $1\leq i\leq n$ and an admissible radius $r$, let $\mu_{i,r}$ be the
uniform probability measure on $B_n(i,r)$, and write
\[
  A_{i,r}z=\sum_{j=1}^n\mu_{i,r}(j)z(j).
\]
Then $(\mathcal A_nz)(i)=\max_r A_{i,r}z$.

\section{The chamber certificate lemma}

Fix signs $s=(s_1,\ldots,s_{n-1})\in\{-1,1\}^{n-1}$ and define
\begin{equation}\label{eq:c-def}
  c_1=s_1,\qquad c_i=s_i-s_{i-1}\ (2\leq i<n),
  \qquad c_n=-s_{n-1}.
\end{equation}
Thus $\sum_i c_i=0$.  The nonzero $c_i$ alternate in sign; internal nonzero
coefficients are $\pm2$, while endpoint coefficients are $\pm1$.

\begin{lemma}[Signed edge-flux certificate]\label{lem:certificate}
Let $z:V_n\to\mathbb R$ and put $y=\mathcal A_nz$.  Fix $s$ and the
coefficients~\eqref{eq:c-def}.  For every $i$ with $c_i>0$, choose a radius
$r_i$ satisfying $y(i)=A_{i,r_i}z$.  Suppose that, for every $i$ with
$c_i<0$, there are nonnegative rational weights $\lambda_{i,r}$ such that
\begin{equation}\label{eq:mass-constraint}
  \sum_r\lambda_{i,r}=-c_i.
\end{equation}
Define the signed measure
\begin{equation}\label{eq:nu-def}
  \nu=
  \sum_{c_i>0}c_i\mu_{i,r_i}
  -\sum_{c_i<0}\sum_r\lambda_{i,r}\mu_{i,r}
\end{equation}
and its edge fluxes
\[
  \Gamma_k=\sum_{j=1}^k\nu(j),\qquad 1\leq k<n.
\]
If $|\Gamma_k|\leq C$ for every $k$, then
\begin{equation}\label{eq:chamber-bound}
  \sum_{i=1}^n c_i y(i)\leq C\Var(z).
\end{equation}
\end{lemma}

\begin{proof}
For a positive coefficient the chosen active radius gives equality.  If
$c_i=-w_i<0$, then $A_{i,r}z\leq y(i)$ for every $r$, so
\[
  c_i y(i)=-w_i y(i)
  \leq-\sum_r\lambda_{i,r}A_{i,r}z
\]
by~\eqref{eq:mass-constraint}.  Summing these estimates yields
\begin{equation}\label{eq:nu-majorant}
  \sum_i c_i y(i)\leq\sum_j\nu(j)z(j).
\end{equation}
The measure $\nu$ has total mass zero because $\sum_i c_i=0$ and every
$\mu_{i,r}$ is a probability measure.  With
$d_k=z(k+1)-z(k)$, Abel summation gives
\[
  \sum_{j=1}^n\nu(j)z(j)
  =-\sum_{k=1}^{n-1}\Gamma_kd_k.
\]
Combining this identity with~\eqref{eq:nu-majorant} proves
\[
  \sum_i c_i y(i)
  \leq\sum_k|\Gamma_k|\,|d_k|
  \leq C\Var(z).
\]
\end{proof}

\begin{corollary}\label{cor:finite-reduction}
Suppose that for every sign vector $s$ and every choice of one radius at each
vertex with $c_i(s)>0$, weights satisfying Lemma~\ref{lem:certificate} exist
with a common bound $C$.  Then $\Cpath n\leq C$.
\end{corollary}

\begin{proof}
For any real sequence $y$,
\[
  \Var(y)=\max_{s\in\{-1,1\}^{n-1}}
  \sum_{k=1}^{n-1}s_k\bigl(y(k)-y(k+1)\bigr)
  =\max_s\sum_i c_i(s)y(i).
\]
For a fixed $z$ and $s$, select any maximizing radius at each positive
coefficient.  This choice occurs in the assumed finite family, including
when active radii are tied.  Lemma~\ref{lem:certificate}, followed by
Lemma~\ref{lem:quotient}, proves the claim.
\end{proof}

\section{The exact certificate for \texorpdfstring{$P_{11}$}{P11}}

At the eleven vertices the numbers of distinct centered balls are
\begin{equation}\label{eq:radius-counts}
  (R_1,\ldots,R_{11})=(11,10,9,8,7,6,7,8,9,10,11).
\end{equation}
We order sign vectors lexicographically with $-1$ before $+1$, and for each
sign vector order the positive active radii lexicographically from radius
zero.  The number of chambers is therefore the independently reproducible
integer
\begin{equation}\label{eq:chamber-count}
  \sum_{s\in\{-1,1\}^{10}}
  \prod_{\{i:c_i(s)>0\}}R_i
  =5{,}874{,}849.
\end{equation}

Most chamber certificates have an especially small form.  List the nonzero
coefficients $c_i$ from left to right.  Each consecutive pair contains one
positive and one negative coefficient.  For each such pair, the certificate
chooses one centered ball at the negative vertex and assigns it weight one.
An internal negative coefficient $-2$ belongs to two consecutive pairs, and
a negative endpoint coefficient $-1$ belongs to one, so the mass
constraint~\eqref{eq:mass-constraint} holds automatically.  The serialized
certificate contains one radius byte per pair.  Across all chambers there
are
\begin{equation}\label{eq:slot-count}
  45{,}011{,}088
\end{equation}
such radius choices.

Ten chambers use rational mixtures instead.  For completeness,
Table~\ref{tab:exceptions} lists all of them.  Chamber indices are zero-based
in the canonical order above; vertices are one-based.  The notation
$v:r[w]$ assigns weight $w$ to the centered ball of radius $r$ at vertex $v$.

\begin{longtable}{@{}r>{\raggedright\arraybackslash}p{0.78\textwidth}@{}}
\caption{The ten rational-mixture exceptions.}\label{tab:exceptions}\\
\toprule
Index & Nonzero negative mixtures \\
\midrule
\endfirsthead
\toprule
Index & Nonzero negative mixtures \\
\midrule
\endhead
50143 & $1:3[44/45]+7[1/45]$; $5:1[2]$; $7:1[2]$; $11:3[44/45]+7[1/45]$.\\
716190 & $1:7[1]$; $3:1[11/10]+5[9/10]$; $5:1[2]$; $7:1[2]$; $9:1[11/10]+5[9/10]$.\\
756109 & As at 716190, together with $11:7[1]$.\\
763791 & $1:7[1]$; $3:1[17/10]+5[3/10]$; $5:1[2]$; $7:1[2]$; $9:1[17/10]+5[3/10]$; $11:7[1]$.\\
1855862 & $2:2[44/45]+6[46/45]$; $5:1[2]$; $7:1[2]$; $10:2[44/45]+6[46/45]$.\\
1860087 & The same mixtures as at 1855862.\\
1864312 & $2:2[152/135]+6[118/135]$; $5:1[2]$; $7:1[2]$; $10:2[152/135]+6[118/135]$.\\
1872762 & $2:2[664/495]+6[326/495]$; $5:1[2]$; $7:1[2]$; $10:2[664/495]+6[326/495]$.\\
5206386 & $3:1[11/10]+5[9/10]$; $5:1[2]$; $7:1[2]$; $9:1[11/10]+5[9/10]$.\\
5250145 & As at 5206386, together with $11:7[1]$.\\
\bottomrule
\end{longtable}

The complete pure-radius data and the exceptional mixtures are supplied as
ancillary files.  The following proposition records exactly what the
independent verifier checks.

\begin{proposition}[Exact finite verification]\label{prop:verification}
For each of the $5{,}874{,}849$ chambers in~\eqref{eq:chamber-count}, the
ancillary certificate satisfies the mass constraints
\eqref{eq:mass-constraint} and the exact edge-flux inequalities
\[
  |\Gamma_k|\leq\frac{10}{11},\qquad 1\leq k\leq10.
\]
\end{proposition}

\begin{proof}
The verifier independently regenerates the sign vectors by bit masks and the
positive-radius tuples by an integer odometer.  It parses exactly the number
of pair radii prescribed by the nonzero coefficients and rejects an invalid
radius, missing data, or trailing data.  In the ten exceptional chambers it
parses positive rational weights and verifies their sum separately at every
negative vertex.

For every chamber it constructs the signed rational vertex measure
\eqref{eq:nu-def}, verifies that its total mass is zero, and computes each
$\Gamma_k$ directly as a rational prefix sum.  It then compares
$|\Gamma_k|$ with $10/11$ by exact integer arithmetic.  The completed run
regenerated all $5{,}874{,}849$ chambers, consumed all
$45{,}011{,}088$ radius slots, verified all ten rational exceptions, and
accepted every edge inequality.  The same verification was repeated from a
fresh temporary directory containing only the verifier and the three
certificate files.
\end{proof}

\begin{proof}[Proof of Theorem~\ref{thm:main}]
Proposition~\ref{prop:verification} supplies the hypotheses of
Corollary~\ref{cor:finite-reduction} with $n=11$ and $C=10/11$.  Hence
$\Cpath{11}\leq10/11$.

For the reverse inequality, let $e_1$ be one at vertex $1$ and zero
elsewhere.  Its variation is one.  At vertex $i$, the smallest centered ball
containing vertex $1$ has radius $i-1$ and cardinality
$\min\{11,2i-1\}$; every larger such ball has at least that cardinality.
Consequently,
\[
  (M_{P_{11}}e_1)(i)=\frac1{\min\{11,2i-1\}}.
\]
This sequence is nonincreasing from $1$ to $1/11$, so its variation is
$10/11$.  Thus $\Cpath{11}\geq10/11$, completing the proof.
\end{proof}

\section{Reproducibility and limitations}

The decisive verification requires Python~3 and only its standard library.
From the root of the ancillary archive, it is run by
\begin{verbatim}
python3 code/verify_radius_certificate.py \
  certificates/p11_radius_certificate.json \
  certificates/p11_radius_certificate.bin
\end{verbatim}
The exceptional-mixture file named by the certificate header is read from the
same directory.  The verifier fails closed on malformed rational data,
invalid radii, incorrect mixture masses, incomplete or surplus payload data,
and inconsistent file bindings.  Deliberate corruptions of each of these
types were rejected.

The certificate generator used linear programming only to locate supports
for the ten rational exceptions.  The displayed weights and every decisive
inequality were subsequently recomputed exactly, and the independent
verifier does not import the generator or any numerical optimization
library.

Theorem~\ref{thm:main} is a finite result.  The enumeration does not provide a
uniform rule for all $n$, so it neither proves the full path conjecture nor
the sharp centered variation inequality on the integer lattice.  The point
masses at the two endpoints show attainment, but we do not classify all
extremizers.

\section*{Data, code, and computational disclosure}

The exact certificate, generator, independent verifier, corruption tests,
and reproduction instructions accompany this manuscript.  Computational
and AI-assisted tools were used for proof discovery, code development,
literature checking, and manuscript preparation.  The mathematical endpoint
depends on the exact rational certificate and the independently executable
verifier described above.  No proof assistant was used.

\bibliographystyle{amsplain}
\bibliography{references}

\end{document}
